\chapter{Results}
\label{ch:results}

This chapter presents the empirical results of the statistical benchmarking suite, evaluating MalthusJAX across the three core hypotheses formulated in Chapter \ref{ch:methodology}: Mathematical Parity (H1), Structural Ablation (H2), and Numerical Representation (H3). The experiments systematically isolate the computational overhead and scaling behaviors of the framework as problem dimensionality $D$ scales from 10 to 1,000.

% ==============================================================================
% 5.1 MATHEMATICAL PARITY (H1)
% ==============================================================================
\section{Mathematical Parity (H1)}
\label{sec:results_parity}

The first requirement of MalthusJAX is to serve as a mathematically identical drop-in replacement for the native Python framework, EvoSAX. To test this, the \texttt{malthusjax\_wrapper} pipeline was executed against the \texttt{evosax\_baseline} across the BBOB testbed (Sphere, Rosenbrock, Rastrigin).

\subsection{Non-Parametric Location Shifts}
Because execution times and fitness landscapes are heavily skewed at varying dimensionalities, we utilize the non-parametric Wilcoxon Signed-Rank test to evaluate location shifts between the paired solvers. Table \ref{tab:wilcoxon_parity} presents the summary statistics and computed $p$-values.

\begin{table}[h]
    \centering
    \caption{Wilcoxon Parity Results for Baseline Equivalence.}
    \label{tab:wilcoxon_parity}
    \resizebox{\textwidth}{!}{
        % INJECT PARITY TABLE ARTIFACT (To be updated with TOST equivalence and Holm-Bonferroni correction)
        \input{results/h1_parity/ols_regression_table.tex}
    }
\end{table}

The computed Cohen's $d_z$ effect sizes are infinitesimally small (e.g., $-0.04$ for the Sphere execution time), indicating that the structural behavior is statistically indistinguishable.

\subsection{Log-Log Scaling Laws}
To guarantee that the JAX JIT compiler overhead does not introduce a hidden asymptotic complexity penalty, we constructed a Log-Log Interaction Regression model:
\begin{equation}
    \ln(\text{Metric}) = \beta_0 + \beta_1 I_{mjx} + \beta_2 \ln(D) + \beta_3 (I_{mjx} \times \ln(D))
\end{equation}

Table \ref{tab:ols_parity} proves that the interaction coefficient ($\beta_3$) is universally non-significant across all execution time scaling regressions (e.g., $p=0.95$ for Rastrigin, $p=0.81$ for Sphere). Consequently, we reject the presence of any asymptotic deviation; MalthusJAX scales identically to EvoSAX.

\begin{table}[h]
    \centering
    \caption{Log-Log OLS Scaling Regression Coefficients.}
    \label{tab:ols_parity}
    \resizebox{\textwidth}{!}{
        % INJECT OLS TABLE ARTIFACT (To be updated with HC0/HC1/HC3 robust standard errors and Holm-Bonferroni correction)
        \input{results/h1_parity/ols_regression_table.tex}
    }
\end{table}

\begin{figure}[h]
    \centering
    \includegraphics[width=0.8\textwidth]{results/h1_parity/malthusjax_wrapper_vs_evosax_baseline_sphere_execution_time_scaling.png}
    \caption{Log-Log scaling of Execution Time vs Dimensionality ($D$) on the Sphere function.}
    \label{fig:parity_sphere_scaling}
\end{figure}

% ==============================================================================
% 5.2 STRUCTURAL ABLATION (H2)
% ==============================================================================
\section{Structural Ablation (H2)}
\label{sec:results_ablation}

Having established baseline parity, we now dissect the specific computational overhead introduced by individual operator compilation. Table \ref{tab:ols_ablation} highlights the differential cost of substituting the wrapped EvoSAX routines with pure JAX-native MalthusJAX operators (Sparse Gaussian Mutation, Uniform Real Crossover, Tournament Selection, Elite Pool).

\begin{table}[h]
    \centering
    \caption{Ablation OLS Regression (Native vs Wrapper).}
    \label{tab:ols_ablation}
    \resizebox{\textwidth}{!}{
        % INJECT ABLATION TABLE ARTIFACT (Updated with HC3 and Holm-Bonferroni)
        \input{results/h2_ablation/ols_regression_table.tex}
    }
\end{table}

\begin{figure}[h]
    \centering
    \includegraphics[width=0.8\textwidth]{results/h2_ablation/mjx_ablate_mutation_vs_mjx_baseline_sphere_execution_time_scaling.png}
    \caption{Scaling impact of native Sparse Gaussian Mutation.}
    \label{fig:ablation_mutation}
\end{figure}

% ==============================================================================
% 5.3 NUMERICAL REPRESENTATION (H3)
% ==============================================================================
\section{Numerical Representation (H3)}
\label{sec:results_representation}

Finally, we evaluate the system's ability to maintain solver convergence fidelity while aggressively down-casting tensor precision. The execution times and best fitness deviations for \texttt{float32}, \texttt{bfloat16}, and \texttt{float16} are benchmarked.

\begin{table}[h]
    \centering
    \caption{Representation Parity (Mixed Precision).}
    \label{tab:wilcoxon_representation}
    \resizebox{\textwidth}{!}{
        % INJECT REPRESENTATION TABLE ARTIFACT (Updated with HC3 and Holm-Bonferroni)
        \input{results/h3_representation/ols_regression_table.tex}
    }
\end{table}

\begin{figure}[h]
    \centering
    \includegraphics[width=0.8\textwidth]{results/h3_representation/mjx_bf16_vs_mjx_f32_sphere_boxplots.png}
    \caption{Location shifts comparing f32 and bf16 on Rosenbrock.}
    \label{fig:representation_bf16_box}
\end{figure}

% ==============================================================================
% 5.4 STRESS TESTING ON HIGHLY NON-LINEAR LANDSCAPES (HARD MODE)
% ==============================================================================
\section{Stress Testing on Highly Non-Linear Landscapes (Hard Mode)}
\label{sec:results_hard_mode}

While the initial benchmarking suite (Sections \ref{sec:results_parity}--\ref{sec:results_representation}) successfully isolated the compilation overhead using foundational testbed functions (Sphere, Rosenbrock, Rastrigin), these topographies are relatively benign and computationally trivial. To rigorously validate the numerical stability of the MalthusJAX architecture—particularly the low-precision convergence capabilities demonstrated in H3—we evaluated the system against the most notoriously difficult, multi-modal, and ill-conditioned landscapes from the Black-Box Optimization Benchmarking (BBOB) suite.

Specifically, the "Hard Mode" suite encompasses:
\begin{itemize}
    \item \textbf{Lunacek (f24)}: Highly multi-modal with deceptive global structure.
    \item \textbf{Schwefel (f20)}: Extreme multi-modality where the global optimum is far from the next best local optima.
    \item \textbf{Gallagher's Gaussian 21-hi (f21)}: 21 optima with severe ill-conditioning and asymmetric peaks.
\end{itemize}

\begin{table}[h]
    \centering
    \caption{Hard Mode Scaling Regression (BBOB f20, f21, f24).}
    \label{tab:hard_mode_regression}
    \resizebox{\textwidth}{!}{
        % INJECT HARD MODE TABLE ARTIFACT
        \input{results/h3_representation_hard/ols_regression_table.tex}
    }
\end{table}

\begin{figure}[h]
    \centering
    \includegraphics[width=0.8\textwidth]{results/h3_representation_hard/mjx_f16_vs_mjx_f32_gallagher_21_hi_best_fitness_scaling.png}
    \caption{Precision Scaling under Extreme Ill-Conditioning (Gallagher 21-hi).}
    \label{fig:hard_mode_gallagher}
\end{figure}

The robust regression results (Table \ref{tab:hard_mode_regression}) confirm the theoretical resilience of the JAX compilation pipeline. Even under extreme topographical stress, the interaction coefficients ($\beta_3$) for low-precision execution remain mathematically indistinguishable from the baseline.
